NUCA architectures introduce performance and security trade-offs that motivate the design of \scheme{}.
This section provides the necessary background on these trade-offs and surveys related work.

\subsection{NUCA Architecture}
Modern multicore processors employ increasingly large cache hierarchies to bridge the performance gap between CPU cores and main memory~\cite{muralimanohar2007interconnect, jaleel2006last}.
As core counts and bandwidth demands increase, last-level caches (LLCs) are often physically distributed into multiple banks or slices rather than implemented as a single monolithic structure.
These LLC banks are connected through an on-chip interconnect or network-on-chip (NoC), which allows multiple cores to access different cache banks concurrently.
This physical distribution gives rise to Non-Uniform Cache Access (NUCA) architectures, where the latency of an LLC access depends on the physical and network distance between the requesting core and the target LLC bank~\cite{kim2002adaptive}.
For example, an access to a nearby LLC bank may traverse only a few NoC hops, while an access to a distant bank may traverse many more.
Consequently, even when two accesses both hit in the LLC, they may exhibit different access latencies due solely to their locations in the on-chip network.

On-chip interconnects can be organized using several topologies, including crossbars, rings, meshes, tori, and 3D tori~\cite{jerger2009chip}.
A typical NoC consists of nodes connected by physical links.
Each node may contain one or more CPU cores, private caches, shared LLC banks, a router, a network interface, coherence-directory state, or memory-controller logic.
A typical NoC router consists of input buffers, a crossbar switch, routing logic, and flow-control logic.
Packets traverse the network hop-by-hop, spending at least one cycle per router for arbitration and switching.
Communication between nodes occurs through packets routed according to the network routing protocol and regulated by flow-control mechanisms.

Among commercial processors, Intel Nehalem used a ring-based interconnect for Xeon processors with up to eight cores, while later Intel Skylake-SP and Skylake-X processors adopted mesh interconnects~\cite{kommrusch2021optimizing,dont-mess-around}.
In distributed LLC designs, an address is typically mapped to an LLC slice and associated coherence metadata.
In Intel manycore systems, this metadata is maintained by Caching and Homing Agents (CHAs), which track the location, ownership, and coherence state of cache lines.
Since CHAs and LLC banks are physically distributed across the chip, accesses to different CHAs or LLC banks can incur different NoC traversal latencies.
This non-uniformity is beneficial for scalability and performance, but it also creates a timing signal that can be exploited by NUCA side-channel attacks.

\subsection{Cache Side-Channel Attacks}
\label{sec:back_sidechannel}
Hardware side-channel attacks exploit observable microarchitectural effects to infer secret-dependent behavior.
Examples of such observable effects include cache hits and misses, access latency in TLBs, caches, and DRAM, and contention in shared hardware structures.
In a typical side-channel attack, the victim performs secret-dependent operations that modify the state or timing behavior of a shared resource, and the attacker observes these changes to infer information about the secret of the victim.

Cache side-channel attacks are among the most widely studied hardware side channels.
They are commonly categorized based on how the attacker prepares the cache state and how the attacker measures the victim-induced change.

\textbf{Prime+Probe~\cite{liu2015last,schwarz2017malware, brasser2017software,kayaalp2016high,aciiccmez2008vulnerability,tromer2010efficient,osvik2006cache,neve2006advances}:}
In Prime+Probe, the attacker first fills selected cache sets with its own cache lines during the prime phase.
After the victim executes, the attacker probes the same sets and measures whether its cache lines were evicted.
Evictions reveal that the victim accessed data mapping to the same cache sets, allowing the attacker to infer secret-dependent memory behavior.

\textbf{Flush+Reload~\cite{gullasch2011cache,gruss2015cache,zhang2014cross,yarom2014recovering,van2015just,irazoqui2014wait,yarom2014flush+}:}
Flush+Reload uses the \texttt{CLFLUSH} instruction to evict a shared cache line and later reloads the same line to determine whether the victim accessed it.
A fast reload indicates that the victim brought the line back into the cache.
This attack provides high spatial and temporal resolution, but it requires the attacker and victim to share memory at cache-line granularity.

\textbf{Flush+Flush~\cite{gruss2016flush+}:}
Flush+Flush exploits the execution-time difference of the \texttt{CLFLUSH} instruction itself.
Because flushing a cached line can take a different amount of time than flushing an uncached line, the attacker can infer whether the victim accessed the target line without explicitly reloading it.
This makes Flush+Flush stealthier than attacks that rely on measuring reload latency.

Recent microarchitectural attacks further show that timing leakage extends beyond traditional cache-set contention; these works are complementary to \scheme{} and target different leakage sources.
For example, HertzBleed~\cite{hertzbleed} shows that dynamic voltage and frequency scaling can turn data-dependent power behavior into timing leakage.
GoFetch~\cite{gofetch} shows that hardware prefetching can undermine constant-time cryptographic implementations.
SQUIP~\cite{squip} exploits contention in shared microarchitectural queues.
Recent defenses such as ClepsydraCache~\cite{clepsydra}, randomized caches, and Goldilocks~\cite{goldilocks} address specific leakage sources such as cache eviction behavior, randomized cache mapping, or frequency-induced timing variation.

\subsection{Existing Defenses and Their Limitations} \label{sec:back_defenses}
Several classes of defenses have been proposed against cache side-channel attacks.
Timing obfuscation techniques perturb or reduce the precision of timing measurements, making it harder for the attacker to decode secret-dependent behavior~\cite{zhang2011predictive,godfrey2013server}.
Partitioning and isolation mechanisms prevent attackers and victims from sharing the same cache
resources~\cite{kiriansky2018dawg,liu2016catalyst,wang2007new,kim2012stealthmem}.
Oblivious RAM (ORAM)~\cite{oram,phantom,sanctum} can mitigate a broad range of access-pattern side channels by hiding memory access patterns, but it generally incurs significant performance and storage overhead due to memory reshuffling and metadata management.
For example, Phantom~\cite{phantom} has $6\times$ performance overhead compared to the baseline system.

Prior defenses against NoC-based side-channel attacks generally attempt to segregate, schedule, or distribute network traffic.
SurfNoc~\cite{wassel2013surfnoc} separates packets based on their application domains and forwards packets from the same domain at the same time.
Reinbrecht et al.~\cite{reinbrecht2016gossip} proposed GossipNoC, which dynamically changes the routing algorithm when abnormal network behavior is detected.
Kar et al.~\cite{remapping-llc} proposed dynamic remapping of LLC banks to defend against contention-based cache side-channel attacks.
Such defenses can reduce interference-based channels because those attacks depend on shared NoC resources or traffic contention.
However, they do not eliminate the fundamental non-uniformity of LLC access latency in NUCA architectures.
In particular, even if contention is removed, a victim access to a nearby LLC bank or CHA can still complete faster than an access to a distant bank or CHA.
This distance-dependent timing signal motivates the need for a defense that equalizes LLC-hit latency rather than only reducing network interference.

\subsection{NoC-Based Side-Channel Attacks}
NoC-based side channels exploit timing or throughput variations in the on-chip interconnect.
In Multi-Processor System-on-Chips (MPSoCs), Reinbrecht \textit{et al.} showed that NoC behavior can be leveraged for Prime+Probe-style cache side-channel attacks~\cite{reinbrecht2016side}.
In this attack, the attacker monitors throughput changes to detect when a cryptographic victim accesses cache lines for key-dependent lookups, and then uses this information to initiate the probe phase.

More recent attacks have demonstrated that timing channels also exist in commercial CPU interconnects.
Paccagnella \textit{et al.} developed timing side-channel attacks on CPU ring interconnects by exploiting NoC and cache contention~\cite{paccagnella2021lotr}.
Dai \textit{et al.} showed that contention in a 2D mesh interconnect of a NUCA system can create measurable timing differences, enabling both covert and side-channel attacks~\cite{dont-mess-around}.
In their attack, the transmitter creates traffic contention on shared NoC paths to encode information, while the attacker detects the transmitted value by measuring access latency.
Similarly, MeshUp exploits mesh-interconnect behavior to construct a stateless cache side channel on modern CPU meshes~\cite{meshup}.

Distance-based NUCA attacks do not rely solely on cache-set conflicts, replacement policies, or eviction behavior.
Instead, they exploit observable latency differences among LLC-hit accesses to different LLC banks or CHAs, where the latency variation arises from NoC distance, congestion, or both.
\scheme{} specifically targets this class of leakage by equalizing the latency of security-sensitive LLC-hit accesses in distributed LLC architectures using delay and/or router-bypass mechanisms.

\subsection{Comparison of Cache Side-Channel Attack Models}
\begin{figure}[t]
    \centering
    \includegraphics[width=\columnwidth]{figs/attack-comparison.pdf}
    \caption{Comparison of communication models between the transmitter (victim) and receiver (attacker) in three cache side-channel attacks: (a) Prime+Probe~\cite{percival2005cache}, (b) Dai et al.~\cite{dont-mess-around}, and (c) the distance-based NUCA attack of     AoK~\cite{mahmud2022distancebased}.}
    \label{fig:attack-comparison}
\end{figure}

\begin{table}[t]
    \centering
    \scriptsize
    \setlength{\tabcolsep}{2.5pt}
    \renewcommand{\arraystretch}{1.18}
    \begin{tabularx}{\columnwidth}{|P{0.58in}|Y|Y|Y|}
        \hline
        \textbf{Property}
        & \makecell[c]{\textbf{Prime+Probe}~\cite{percival2005cache}}
        & \makecell[c]{\textbf{Dai et al.}~\cite{dont-mess-around}}
        & \makecell[c]{\textbf{AoK}~\cite{mahmud2022distancebased}} \\
        \hline \hline
        Target arch.
        & Any cache or NUCA
        & Ring-based mesh
        & Multi-hop mesh \\
        \hline
        Covert?
        & No
        & Yes
        & No \\
        \hline
        Attack model
        & Eviction-based contention
        & Traffic-based contention
        & Distance-based timing \\
        \hline
        Exploit target
        & Shared cache sets
        & Shared NoC paths/rings
        & LLC bank/CHA distance \\
        \hline
        Defenses
        & ORAM; partitioning
        & ORAM; clustering; scheduling
        & ORAM; delay/bypass \\
        \hline
    \end{tabularx}
    \caption{Comparison of representative cache and NoC side-channel attacks and their applicable defense categories.}
    \label{tab:attack_differences}
\end{table}

We now compare the communication models of three representative cache side-channel attacks: Prime+Probe~\cite{percival2005cache}, the NoC-contention attack by Dai \textit{et al.}~\cite{dont-mess-around}, and the distance-based NUCA attack demonstrated by Attack of Knights (AoK)~\cite{mahmud2022distancebased}.
Figure~\ref{fig:attack-comparison} illustrates the communication relationship between the transmitter (victim) and the receiver (attacker) in these attacks, while Table~\ref{tab:attack_differences} summarizes their key differences.

Prime+Probe is an eviction-based contention channel.
The victim transmits information implicitly by evicting cache lines of the attacker from shared cache sets, and the attacker decodes the information by measuring which of its lines were evicted.
In contrast, the attack by Dai \textit{et al.} is a traffic-based contention channel.
The transmitter generates NoC traffic on paths shared with the receiver, and the receiver detects the transmitted value by measuring the latency increase caused by network contention.

AoK follows a different communication model.
Rather than relying on eviction-based cache contention or deliberate NoC traffic contention, AoK exploits distance-dependent timing differences in NUCA architectures.
The attacker observes the timing of victim operations and infers whether the victim accessed a near or far LLC bank or CHA.
Therefore, a distance-based NUCA attack requires an architecture in which LLC-hit accesses to different banks or CHAs produce distinguishable timing differences, such as a multi-hop mesh-based manycore processor.

This distinction has important implications for defenses.
Because Prime+Probe and the attack by Dai \textit{et al.} rely on shared resources, they can be mitigated by partitioning, clustering, scheduling, or otherwise reducing resource sharing between the attacker and victim.
However, distance-based NUCA attacks cannot be fully eliminated by partitioning or clustering alone, as long as the attacker can time victim operations.
Their success depends on the observability of latency differences, not necessarily on direct contention over a shared cache set or NoC path.
ORAM~\cite{oram} can mitigate all three classes of attacks by hiding memory access patterns, but it generally incurs high overhead.
\scheme{} therefore takes a more targeted approach: it removes the timing signal exploited by NUCA attacks by enforcing constant latency for security-sensitive LLC-hit accesses, without requiring memory access pattern shuffling or the associated storage overhead.
